% Section 4: The Substrate Question

The previous section revealed that impossibility claims cluster at the implementational
level \cite{campero2025taxonomy}. This section confronts the substrate question directly:
what physical properties, if any, are necessary for consciousness, and can digital systems
possess them?

\subsection{Functionalism vs.\ Substrate Dependence}

The classical functionalist position, articulated most forcefully by Putnam's multiple
realizability argument, holds that mental states are defined by their causal-functional
roles rather than their physical constitution \cite{chalmers2023could}. If consciousness is
a functional property---what it \emph{does} rather than what it \emph{is made of}---then
substrate should not matter. A sufficiently organized system of silicon transistors,
mechanical gears, or even water pipes could in principle realize consciousness.

This view underwrites Camp 1's optimism about AI consciousness. If GWT correctly identifies
consciousness with a broadcast architecture \cite{baars1988cognitive, dehaene2014toward},
and if broadcast architectures can be implemented in software, then software can be conscious.

Integrated Information Theory challenges this picture at its foundation. According to IIT
\cite{tononi2004information, tononi2016integrated, albantakis2023iit4}, consciousness is
identical to integrated information ($\Phi$)---a measure of a system's causal structure.
Critically, $\Phi$ is \emph{not} a functional property in the sense functionalists intend.
Two systems with identical input-output mappings can have radically different $\Phi$ values.

Consider two systems that both implement the NAND gate function:

\begin{itemize}[leftmargin=2em]
\item System A: Two transistors arranged such that the output causally depends on both
      inputs through irreducible interaction. The minimum information partition cuts through
      both inputs simultaneously, yielding non-zero $\Phi$.
\item System B: A lookup table that maps input pairs to outputs. Each entry is causally
      independent. The minimum information partition cuts between individual entries,
      yielding $\Phi \approx 0$.
\end{itemize}

Both systems realize the same function, but only System A possesses integrated information.
IIT thus denies multiple realizability for consciousness: not every physical realization of
a given function instantiates the same causal structure.

The Free Energy Principle \cite{friston2010free, wiese2024fep} extends this critique in a
thermodynamic direction, emphasizing Markov blankets, active inference, and non-equilibrium
steady states as \emph{physical} properties that cannot be captured by functional description.
A digital simulation of a self-organizing system computes the equations governing its dynamics
but does not thereby enter into thermodynamic coupling with an environment.

This distinction between simulation and replication is crucial. A weather model simulates
atmospheric dynamics without becoming weather. A neural network simulation computes state
transitions without becoming a biological network. The functionalist must explain why
consciousness is different---why simulating the right informational structure suffices for
consciousness when it does not suffice for other physical properties.

\subsection{What the Campero Taxonomy Reveals}

The Campero et al.\ taxonomy \cite{campero2025taxonomy} provides empirical evidence that
substrate concerns dominate the impossibility arguments. Of the D3-level objections
(fundamental barriers to AI consciousness), five explicitly invoke implementational
properties:

\begin{enumerate}[leftmargin=2em]
\item \textbf{Electromagnetic field theories}: Consciousness requires specific EM field
      dynamics that arise in biological neurons but not in silicon transistors.
\item \textbf{IIT causal structure}: Digital systems typically lack the dense causal
      integration required for non-zero $\Phi$.
\item \textbf{Analog processing}: Consciousness requires continuous-valued dynamics rather
      than discrete symbolic computation.
\item \textbf{Quantum effects}: Some theories argue consciousness depends on quantum
      coherence in microtubules, which cannot be replicated in classical substrates.
\item \textbf{G\"{o}delian limits}: If consciousness involves non-computable operations,
      then Turing machines cannot replicate it by definition.
\end{enumerate}

Meanwhile, algorithmic-level concerns are predominantly D1 or D2 objections. The asymmetry
is clear: \textbf{skepticism about AI consciousness is skepticism about digital substrates,
not algorithms.}

This has profound methodological implications. Debates about whether transformers have the
``right'' attention mechanism or whether reinforcement learning agents have the ``right''
temporal integration are largely orthogonal to the core question. Even an AI system with
perfect functional alignment to biological consciousness could still fail to be conscious if
substrate matters.

Conversely, substrate-dependent theories face a measurement problem. IIT requires computing
the minimum information partition to determine $\Phi$, which is NP-hard and intractable for
large systems \cite{oizumi2014phenomenology}. FEP requires identifying Markov blankets and
measuring thermodynamic flows, which presupposes knowing where the system's boundaries are.
For complex AI systems, these measurements are currently impossible, leaving substrate-dependent
theories empirically idle.

\subsection{The Brain Organoid Stress Test}

Brain organoids---three-dimensional cultures of neural tissue derived from stem cells---provide
a crucial test case that exposes tensions within every theoretical camp. Organoids are
increasingly sophisticated: they exhibit spontaneous oscillatory activity, form layered
structures resembling cortical columns, and develop rudimentary connectivity patterns. Yet they
lack sensory input, motor output, and any behavioral manifestation of consciousness.

For \textbf{Integrated Information Theory}, organoids pose a direct challenge. If $\Phi$
depends only on causal structure, and organoids possess genuine biological neural networks
with recurrent connectivity, then they should have non-zero $\Phi$. As organoids grow more
complex---developing synchronized oscillations and distributed information processing---their
$\Phi$ should increase. IIT appears committed to the conclusion that sufficiently complex
organoids possess some degree of micro-consciousness, even without any sensory experience or
behavioral expression. This is not a reductio---IIT explicitly embraces the possibility of
micro-consciousness in simple systems \cite{tononi2016integrated}---but it raises pressing
ethical questions about organoid research.

For \textbf{Global Workspace Theory}, organoids are clearly unconscious. GWT requires a
broadcast architecture connecting specialized modules through a global workspace, typically
identified with thalamocortical loops and frontoparietal networks \cite{baars1988cognitive,
dehaene2014toward}. Organoids lack this architecture entirely. This verdict is revealing:
GWT \emph{denies consciousness to biological tissue}. The theory is not substrate-neutral
in practice---it requires specific neural architectures that happen to exist in biological
brains. But this undermines the functionalist credentials GWT is often assumed to possess.
If biological substrate is insufficient without the right architecture, why should silicon
substrate suffice \emph{with} the right architecture? The asymmetry needs justification.

For \textbf{Free Energy Principle} approaches, organoids occupy ambiguous territory. They are
clearly self-organizing---they develop structure autonomously through local cellular
interactions and exhibit non-equilibrium dynamics. This satisfies part of Wiese's criterion
\cite{wiese2024fep}. However, organoids lack sensorimotor coupling. They do not engage in
active inference because they have no motors to act with and no coherent sensory input to
predict from. Does self-organization without active inference suffice for consciousness under
FEP? The theory does not provide a clear answer.

For \textbf{Sritriratanarak and Garcia's boolean-switch model} \cite{sritriratanarak2025consciousness},
organoids present a different puzzle. Organoids possess reactive substrates---neurons fire,
connections strengthen, electrical activity propagates automatically. But organoids do not
engage in hypothetical reasoning. There is nothing for the consciousness switch to suppress.
Does the switch engage spontaneously? The theory provides no mechanism for deciding.

\subsection{Implications}

The organoid case reveals that every major theory contains hidden commitments about what
``counts'' for consciousness:

\begin{itemize}[leftmargin=2em]
\item IIT: Causal structure alone suffices, even without behavior or experience content.
\item GWT: Biological substrate does not suffice without functional architecture.
\item FEP: Self-organization might not suffice without environmental coupling.
\item Boolean-switch model: Reactive substrates might not suffice without cognitive functions
      to suppress.
\end{itemize}

None of these commitments are obviously wrong, but they are \emph{underspecified}. Each
theory addresses a subset of consciousness's properties while leaving others unaccounted for.
This suggests that the search for a single necessary-and-sufficient condition is misguided.
Consciousness may be a cluster concept \cite{butlin2023consciousness}, requiring multiple
jointly sufficient conditions that can be satisfied to varying degrees.

If so, then the question ``Can AI be conscious?'' has no binary answer. Instead, we should
ask: \emph{Which properties associated with consciousness can AI systems instantiate, and
which are substrate-dependent?} The organoid case serves as a stress test that every theory
fails in instructive ways. A full analysis of organoid consciousness, exploring these tensions
in detail and proposing experimental protocols, is developed in a companion paper.
